\scanpage{316}
\mspage{15}
% archive sheet number: 317

\uline{Remarques} a) On trouve essentiellement
les $\overline{\mathrm{m}}$s objets ($L_{\ast}$ ou $L_{\cdot}$) pour exprimer
des \uline{opérations} sur les invariants additifs
contravariants, ou sur invariants \inserted{``additifs''} covariants
(du type envisagé) -- \unsure{c'était|s'étant}
clair à priori, à cause du passage
de $A$ à $B$ par \inserted{passage à la} catégorie opposée
-- mais l'op. définie par $L_{\ast}$ dépend
de $L_{\ast}$ de façon contravariante dans le
premier cas, covariante dans le deuxième.

b) On a en principe résolu, et de façon
quasi-tautologique, la question initiale
des \deleted{\unsure{ce qu'à faire sur}} \inserted{déterminer les constructions} possibles sur
invariants contravariants ou covariants
\inserted{``additifs''} du type envisagé : On trouve une
\inserted{$\overline{\mathrm m}$} grosse catégorie, $\mathit{Ab}_{\ast}$ ou $\mathit{Ab}_{\cdot}$ (\unsure{au|ou}
passage près de celles-ci à l'opposée). Cela
tient au fait que l'on a admis
des catégories \inserted{$A$} de valeurs \inserted{$B$ ou $A$} assez spéciales,
à savoir où on peut effectuer \uline{toutes}
les $\varprojlim$, ou ttes les $\varinjlim$ -- ce
qui a pour effet que ces $\overline{\mathrm m}$ opérations
peuvent s'effectuer sur les \deleted{opéra\unsure{tions}}
\deleted{foncteurs} invariants, et \unsure{opé|op.} sur les invariants.

On trouvera que lorsqu'on est plus
restrictif sur les opérations qu'on
se permet, \deleted{(p. ex. si} \inserted{donc} moins restrictif
sur $A$ et $B$, on trouve des catégories
correspondantes d'opérations (ou constructions)
sur invariants plus petites -- \deleted{alors}
en $\overline{\mathrm m}$ temps qu'on trouve des cat.
``universelles'' plus petites \deleted{\unsure{où les}} \inserted{pour tous les}
invariants additifs contrav. ou cov.
possibles. \deleted{\unsure{g}} De façon précise, il doit
résulter des principes généraux (que je
n'ai pas détaillés) que si $\deleted{\unsure{\mathbb C}_{\ast}}F_0:\deleted{\unsure{\mathit{Ab}}}\inserted{\Delta^{\circ}\to B_0}$
est l'objet de $\bigl(\underbrace{(\mathit{Ab}_{\ast})^{\circ}}_{B_0}\bigr)_{\ast}
\simeq(\mathit{Ab}_{\ast}^{\deleted{\unsure{\ast}}\ast})^{\circ}$
qui exprime l'objet \deleted{\unsure{$A^{\circ}$}} $1/2$ simplicial
additif universel, alors la \deleted{sous} catégorie
cherchée d'opérations [correspondant à
un paquet de de cat. d'indices par
rapport auxquelles on se propose de
prendre soit des $\varprojlim$ -- dans des cat.
$B$ où on regarde des inv. de nature
contravariante -- soit des $\varinjlim$ -- dans
le cas de cat. $A$ où on regarde les inv. de
nature cov --] est la sous-catégorie \uline{pleine}
\inserted{$\mathcal U_L$} de $B_0=(\mathit{Ab}_{\ast})^{\circ}$ engendrée par \inserted{les} $F_0(\Delta_{\ast}^{n})$ grâce
aux $\varinjlim$ de type $L$, et \deleted{comme} \inserted{correspond}
donc à la ss-catégorie pleine de $\mathit{Ab}_{\ast}$
engendrée par \inserted{les} $G_0(\Delta^n)$ \deleted{par} \inserted{grâce aux} $\varprojlim$ de type $L$.
